% References to labels of other documents with both xr and xr-hyper
% loaded, xr first, then hyperref. The page builds, and each reference
% prints as with one of them: from an .aux that LaTeX wrote under
% hyperref, whose \externaldocument gives the URL of its PDF, its number,
% page, name, and \autoref as in the PDF, with no link and no URL, in each
% context; from an .aux that LaTeX wrote without hyperref, its number and
% page; from an .aux that tex4ht wrote under hyperref, with its .xref, a
% link to that document's page, also from an \externaldocument before
% xr-hyper loads; one through \externalcitedocument, with no prefix; a
% label of an \@input file of the other .aux; and a label of a file that
% does not exist yet, as ??.
\documentclass{article}
\usepackage{amsthm}
\usepackage{xr}
\externaldocument[E-]{test-xr/hyper4ht-hyper}
\usepackage{xr-hyper}
\usepackage{hyperref}
\externaldocument[X-]{test-xr/latex}[https://example.org/ext.pdf]
\externaldocument[H-]{test-xr/hyper4ht-hyper}
\externalcitedocument{test-xrurl/tex4ht}
\externaldocument[P-]{test-xr/plain}
\externaldocument{test-xrboth-none}
\newtheorem{theorem}{Theorem}
\begin{document}
\section{See \ref{X-ss}}
CaseA \ref{X-ss}, page \pageref{X-front}, named \nameref{X-thm},
auto \autoref{X-thm}.

CaseB \ref{H-oss}, auto \autoref{H-oss}.

CaseC \ref{oss}.

CaseEarly \ref{E-oss}.

CaseD \ref{P-ss}, page \pageref{P-front}.

CaseE \ref{none}.

Included \ref{X-c}.

\textit{In italics, \ref{X-ss}.}
A note.\footnote{In a footnote, \ref{X-ss}.}
\textsf{In sans-serif, \ref{X-ss}.}
\texttt{In typewriter, \ref{X-ss}.}

\begin{theorem}
In a theorem, \ref{X-ss}.
\end{theorem}

\section{Here}\label{here}
Local \ref{here}.
\end{document}
